Write \(\psi=\psi_{\eta}\).  The exact Luxemburg premise gives
\[
 E_P\exp(\eta^2/\psi^2)\le2,
 \qquad E_P\eta=0.
\]
We first derive a numerical growth bound.  Let \(\eta'\) be an independent copy of \(\eta\).  If \(|t|\psi\le1\), Jensen's inequality and centering give
\[
 E_Pe^{t\eta}\le E_Pe^{t(\eta-\eta')}.
\]
The difference \(\eta-\eta'\) is symmetric.  Using \(\cosh u\le e^{u^2/2}\), \((\eta-\eta')^2\le2\eta^2+2(\eta')^2\), and then the Lyapunov inequality with exponent \(t^2\psi^2\le1\), we obtain
\[
\begin{aligned}
 E_Pe^{t\eta}
 &\le E_P\cosh\{t(\eta-\eta')\}\\
 &\le E_P\exp\{t^2(\eta-\eta')^2/2\}\\
 &\le \{E_Pe^{t^2\eta^2}\}^2
 \le 2^{2t^2\psi^2}
 \le e^{4\psi^2t^2}.
\end{aligned}
\]
If \(|t|\psi>1\), Young's inequality and Cauchy--Schwarz yield
\[
 t\eta\le\frac{\eta^2}{2\psi^2}+\frac{t^2\psi^2}{2},
 \qquad
 E_Pe^{t\eta}\le\sqrt2\,e^{t^2\psi^2/2}
 \le e^{4\psi^2t^2}.
\]
Consequently, for every \(z\in\mathbb C\),
\[
 |M(z)|\le E_Pe^{\operatorname{Re}(z)\eta}
 \le \exp(4\psi^2|z|^2). \tag{1}
\]

Suppose for contradiction that \(M\) has no zero in the closed disk of radius \(R_0\).  Since \(M(0)=1\), an analytic logarithm \(f=\log M\), normalized by \(f(0)=0\), exists on a neighborhood of that disk.  Centering gives \(f'(0)=0\).  For \(0<R<R_0\), (1) gives
\[
 \sup_{|z|\le R}\operatorname{Re}f(z)\le4\psi^2R^2.
\]
Apply the Carath\'eodory coefficient estimate on the radius-\(R\) disk to the positive-real-part analytic function \(4\psi^2R^2-f\).  Its value at zero is \(4\psi^2R^2\), and hence
\[
 |f^{(k)}(0)|
 \le 2k!R^{-k}\,4\psi^2R^2
 =8k!\psi^2R^{2-k}. \tag{2}
\]
Letting \(R\uparrow R_0\), using \(f^{(k)}(0)=\kappa_k(\eta)\), and invoking \(\text{ass:cumulant-separation}\), we get
\[
 \delta\le8k!\psi^2R_0^{2-k}.
\]
But the displayed definitions of \(A_k\) and \(R_0\) give
\[
 8k!\psi^2R_0^{2-k}
 =\frac{8k!}{2^{k+4}k!}\,\delta
 =2^{-k-1}\delta<\delta,
\]
a contradiction.  Thus \(M\) has a zero in \(\{|z|\le R_0\}\), with the asserted numerical radius.